\documentclass[11pt]{article}
\usepackage[a4paper,margin=1in]{geometry}
\usepackage{amsmath,amssymb,mathtools,bm}
\usepackage{enumitem,booktabs,array}
\usepackage{tcolorbox}
\usepackage{hyperref}
\usepackage[protrusion=true,expansion=false]{microtype}
\hypersetup{colorlinks=true,linkcolor=blue,urlcolor=blue,citecolor=blue}
\setlist[itemize]{topsep=2pt,itemsep=2pt}
\setlist[enumerate]{topsep=2pt,itemsep=2pt}
\newtcolorbox{keybox}{colback=blue!3!white,colframe=blue!40!black,boxrule=0.4pt,arc=1mm}
\newtcolorbox{alertbox}{colback=red!3!white,colframe=red!40!black,boxrule=0.4pt,arc=1mm}
\newtcolorbox{answerbox}{colback=green!3!white,colframe=green!40!black,boxrule=0.4pt,arc=1mm}
\newtcolorbox{drillbox}{colback=yellow!5!white,colframe=orange!60!black,boxrule=0.4pt,arc=1mm}
\newcommand{\EV}{\mathbb{E}}
\newcommand{\Var}{\mathrm{Var}}
\newcommand{\Cov}{\mathrm{Cov}}
\newcommand{\blank}[1]{\underline{\hspace{#1}}}

\title{W11-01: Sharfman \& Fernando (2008, SMJ) \\
\large ``Environmental Risk Management and the Cost of Capital'' --- Active Recall Workbook}
\author{Banking \& Risk Management --- Exam Prep}
\date{\today}
\begin{document}\maketitle

\begin{keybox}
\textbf{Paper in one sentence.} Cross-section of 267 S\&P 500 firms in 2002: firms with better environmental risk management (KLD scores plus EPA TRI emissions) enjoy a lower cost of equity ($\sim$1--2 pp) and higher leverage --- a link between corporate environmental performance and financial outcomes.

\textbf{Placement.} Early, influential ESG-and-cost-of-capital paper. Predates the climate-risk asset-pricing literature; complements Chava (W11-03) on debt cost and Hong-Kacperczyk (W11-02) on sin stocks.
\end{keybox}

\section*{Round 1: Complete Walk-Through}

\subsection*{1.1 Data}
267 S\&P 500 firms, 2002. Environmental performance measured by (i) KLD environmental strengths minus concerns, (ii) EPA Toxics Release Inventory (TRI) emissions. Cost of equity from ValueLine / implied cost-of-capital models. Leverage from Compustat.

\subsection*{1.2 Central regression}
\[
\text{CostEquity}_i = \alpha + \beta\cdot\text{EnvPerformance}_i + X_i'\gamma + \varepsilon_i.
\]
Expected $\beta<0$: better environmental performance lowers cost of equity.

\subsection*{1.3 Headline results}
\begin{itemize}
\item Firms with better environmental performance have $\sim$1--2 pp lower implied cost of equity.
\item Same firms have higher leverage ratios --- they can sustain more debt.
\item Result robust to industry controls.
\end{itemize}

\subsection*{1.4 Mechanism}
\begin{itemize}
\item Better ERM reduces tail-risk exposure (spills, litigation, regulation).
\item Lower perceived risk $\Rightarrow$ lower required return by investors.
\item Lower distress probability $\Rightarrow$ higher debt capacity.
\end{itemize}

\subsection*{1.5 Caveats}
\begin{alertbox}
\begin{itemize}
\item Cross-sectional, one-year snapshot; causality not identified.
\item KLD scores have measurement issues; TRI covers only US plants.
\item Omitted-variable bias: well-governed firms manage environment \emph{and} have lower cost of equity.
\end{itemize}
\end{alertbox}

\section*{Round 2: Level 1 Fill-in}
\begin{enumerate}
\item Sample: \blank{1cm} S\&P 500 firms, year \blank{1cm}.
\item Env performance measures: \blank{1cm} score and EPA \blank{1cm}.
\item Cost-of-equity differential: $\sim$\blank{1cm}--\blank{1cm} pp lower.
\item Leverage direction: \blank{1cm}.
\item Mechanism: lower \blank{1cm}-risk exposure.
\end{enumerate}

\section*{Round 3: Level 2 Prompts}
\begin{enumerate}
\item Write the central spec and discuss identifying assumptions.
\item How does ERM reduce the firm's systematic or idiosyncratic risk?
\item Discuss omitted-variable concerns.
\item Contrast SF with Chava (2014) on debt markets.
\end{enumerate}

\section*{Round 4: Minimal Scaffolding}
From a blank page: (i) data, (ii) measures, (iii) spec, (iv) results, (v) caveats.

\section*{Round 5: Blank Page}
Essay: does environmental performance lower the cost of capital? SF evidence.

\vspace{6pt}
\begin{drillbox}
\textbf{Speed Drill.} (1) Sample size and year. (2) Two env measures. (3) Cost-of-equity delta. (4) Leverage direction. (5) One caveat.
\end{drillbox}

\section*{Quick Reference \& Answer Key}
\begin{answerbox}
\begin{itemize}
\item \textbf{Sample.} 267 S\&P 500, 2002.
\item \textbf{Measures.} KLD scores $+$ EPA TRI.
\item \textbf{Cost of equity.} $\sim$1--2 pp lower for good ERM.
\item \textbf{Leverage.} Higher.
\end{itemize}
\end{answerbox}

\begin{keybox}
\textbf{30-second exam pitch.} Sharfman \& Fernando (2008) establish an early and influential cross-sectional link between corporate environmental performance and financial outcomes. Using 267 S\&P 500 firms in 2002, measuring environmental performance with KLD scores and EPA Toxics Release Inventory emissions, they show that firms with better environmental risk management have $\sim$1--2 pp lower implied cost of equity and higher leverage. The mechanism: strong ERM reduces tail risks from spills, litigation, and regulation, which lowers investor risk premia and permits more debt. The paper is a forerunner of the modern climate-and-ESG asset-pricing literature; causal interpretation is limited by the cross-sectional design and by omitted-variable concerns around corporate governance, but the correlation is strong and has been replicated broadly.
\end{keybox}

\end{document}
